\section{An explicit construction at every height}\label{sec:general}

We now prove Proposition~\ref{thm:general}. We place the doubles far apart and let the
singletons between them be black in a boundary row. Call a singleton a \emph{sea} if its
black point lies in row $1$ or in row $h$.

\begin{proposition}[alternating seas]\label{prop:seas}
Let $h\ge3$ and let $H$ be a triangle-free graph on the rows with edges $e_1,\dots,e_q$,
where $q\ge1$. Let $Z$ be the blackout of the $h\times m$ grid with the following columns.
The columns $p_1<\dots<p_q$ are doubles, where column $p_j$ is black exactly in the 2 rows
of $e_j$ and $p_{j+1}-p_j\ge W_h+1$. Every other column is a sea. The seas to the left of
$p_1$ are black in row $1$, and the row of the seas switches between $1$ and $h$ at every
double. Then $Z$ is valid and $|Z|=m+q$.
\end{proposition}

\begin{figure}[ht]
\centering
\begin{tikzpicture}[x=0.62cm,y=0.62cm,baseline=(current bounding box.north)]
  \begin{scope}[xshift=0.000cm]
    \draw[black,line width=0.25pt,dash pattern=on 0.93pt off 0.93pt] (1,-1) -- (15,-1) (1,-2) -- (15,-2) (1,-3) -- (15,-3) (1,-4) -- (15,-4) (1,-1) -- (1,-4) (2,-1) -- (2,-4) (3,-1) -- (3,-4) (4,-1) -- (4,-4) (5,-1) -- (5,-4) (6,-1) -- (6,-4) (7,-1) -- (7,-4) (8,-1) -- (8,-4) (9,-1) -- (9,-4) (10,-1) -- (10,-4) (11,-1) -- (11,-4) (12,-1) -- (12,-4) (13,-1) -- (13,-4) (14,-1) -- (14,-4) (15,-1) -- (15,-4);
    \filldraw[fill=kgray,draw=black,line width=0.37pt] (1,-1) circle (0.3);
    \draw[black,line width=0.37pt] (0.788,-0.788) -- (1.212,-1.212) (0.788,-1.212) -- (1.212,-0.788);
    \filldraw[fill=white,draw=black,line width=0.37pt] (2,-1) circle (0.3);
    \filldraw[fill=white,draw=black,line width=0.37pt] (3,-1) circle (0.3);
    \filldraw[fill=white,draw=black,line width=0.37pt] (4,-1) circle (0.3);
    \filldraw[fill=kgray,draw=black,line width=0.37pt] (5,-1) circle (0.3);
    \draw[black,line width=0.37pt] (4.788,-0.788) -- (5.212,-1.212) (4.788,-1.212) -- (5.212,-0.788);
    \filldraw[fill=kgray,draw=black,line width=0.37pt] (6,-1) circle (0.3);
    \draw[black,line width=0.37pt] (5.788,-0.788) -- (6.212,-1.212) (5.788,-1.212) -- (6.212,-0.788);
    \filldraw[fill=kgray,draw=black,line width=0.37pt] (7,-1) circle (0.3);
    \draw[black,line width=0.37pt] (6.788,-0.788) -- (7.212,-1.212) (6.788,-1.212) -- (7.212,-0.788);
    \filldraw[fill=kgray,draw=black,line width=0.37pt] (8,-1) circle (0.3);
    \draw[black,line width=0.37pt] (7.788,-0.788) -- (8.212,-1.212) (7.788,-1.212) -- (8.212,-0.788);
    \filldraw[fill=white,draw=black,line width=0.37pt] (9,-1) circle (0.3);
    \filldraw[fill=white,draw=black,line width=0.37pt] (10,-1) circle (0.3);
    \filldraw[fill=white,draw=black,line width=0.37pt] (11,-1) circle (0.3);
    \filldraw[fill=white,draw=black,line width=0.37pt] (12,-1) circle (0.3);
    \filldraw[fill=white,draw=black,line width=0.37pt] (13,-1) circle (0.3);
    \filldraw[fill=kgray,draw=black,line width=0.37pt] (14,-1) circle (0.3);
    \draw[black,line width=0.37pt] (13.788,-0.788) -- (14.212,-1.212) (13.788,-1.212) -- (14.212,-0.788);
    \filldraw[fill=kgray,draw=black,line width=0.37pt] (15,-1) circle (0.3);
    \draw[black,line width=0.37pt] (14.788,-0.788) -- (15.212,-1.212) (14.788,-1.212) -- (15.212,-0.788);
    \filldraw[fill=white,draw=black,line width=0.37pt] (1,-2) circle (0.3);
    \filldraw[fill=white,draw=black,line width=0.37pt] (2,-2) circle (0.3);
    \filldraw[fill=white,draw=black,line width=0.37pt] (3,-2) circle (0.3);
    \filldraw[fill=white,draw=black,line width=0.37pt] (4,-2) circle (0.3);
    \filldraw[fill=white,draw=black,line width=0.37pt] (5,-2) circle (0.3);
    \filldraw[fill=white,draw=black,line width=0.37pt] (6,-2) circle (0.3);
    \filldraw[fill=white,draw=black,line width=0.37pt] (7,-2) circle (0.3);
    \filldraw[fill=white,draw=black,line width=0.37pt] (8,-2) circle (0.3);
    \filldraw[fill=kgray,draw=black,line width=0.37pt] (9,-2) circle (0.3);
    \draw[black,line width=0.37pt] (8.788,-1.788) -- (9.212,-2.212) (8.788,-2.212) -- (9.212,-1.788);
    \filldraw[fill=white,draw=black,line width=0.37pt] (10,-2) circle (0.3);
    \filldraw[fill=white,draw=black,line width=0.37pt] (11,-2) circle (0.3);
    \filldraw[fill=white,draw=black,line width=0.37pt] (12,-2) circle (0.3);
    \filldraw[fill=kgray,draw=black,line width=0.37pt] (13,-2) circle (0.3);
    \draw[black,line width=0.37pt] (12.788,-1.788) -- (13.212,-2.212) (12.788,-2.212) -- (13.212,-1.788);
    \filldraw[fill=white,draw=black,line width=0.37pt] (14,-2) circle (0.3);
    \filldraw[fill=white,draw=black,line width=0.37pt] (15,-2) circle (0.3);
    \filldraw[fill=kgray,draw=black,line width=0.37pt] (1,-3) circle (0.3);
    \draw[black,line width=0.37pt] (0.788,-2.788) -- (1.212,-3.212) (0.788,-3.212) -- (1.212,-2.788);
    \filldraw[fill=white,draw=black,line width=0.37pt] (2,-3) circle (0.3);
    \filldraw[fill=white,draw=black,line width=0.37pt] (3,-3) circle (0.3);
    \filldraw[fill=white,draw=black,line width=0.37pt] (4,-3) circle (0.3);
    \filldraw[fill=white,draw=black,line width=0.37pt] (5,-3) circle (0.3);
    \filldraw[fill=white,draw=black,line width=0.37pt] (6,-3) circle (0.3);
    \filldraw[fill=white,draw=black,line width=0.37pt] (7,-3) circle (0.3);
    \filldraw[fill=white,draw=black,line width=0.37pt] (8,-3) circle (0.3);
    \filldraw[fill=kgray,draw=black,line width=0.37pt] (9,-3) circle (0.3);
    \draw[black,line width=0.37pt] (8.788,-2.788) -- (9.212,-3.212) (8.788,-3.212) -- (9.212,-2.788);
    \filldraw[fill=white,draw=black,line width=0.37pt] (10,-3) circle (0.3);
    \filldraw[fill=white,draw=black,line width=0.37pt] (11,-3) circle (0.3);
    \filldraw[fill=white,draw=black,line width=0.37pt] (12,-3) circle (0.3);
    \filldraw[fill=white,draw=black,line width=0.37pt] (13,-3) circle (0.3);
    \filldraw[fill=white,draw=black,line width=0.37pt] (14,-3) circle (0.3);
    \filldraw[fill=white,draw=black,line width=0.37pt] (15,-3) circle (0.3);
    \filldraw[fill=white,draw=black,line width=0.37pt] (1,-4) circle (0.3);
    \filldraw[fill=kgray,draw=black,line width=0.37pt] (2,-4) circle (0.3);
    \draw[black,line width=0.37pt] (1.788,-3.788) -- (2.212,-4.212) (1.788,-4.212) -- (2.212,-3.788);
    \filldraw[fill=kgray,draw=black,line width=0.37pt] (3,-4) circle (0.3);
    \draw[black,line width=0.37pt] (2.788,-3.788) -- (3.212,-4.212) (2.788,-4.212) -- (3.212,-3.788);
    \filldraw[fill=kgray,draw=black,line width=0.37pt] (4,-4) circle (0.3);
    \draw[black,line width=0.37pt] (3.788,-3.788) -- (4.212,-4.212) (3.788,-4.212) -- (4.212,-3.788);
    \filldraw[fill=kgray,draw=black,line width=0.37pt] (5,-4) circle (0.3);
    \draw[black,line width=0.37pt] (4.788,-3.788) -- (5.212,-4.212) (4.788,-4.212) -- (5.212,-3.788);
    \filldraw[fill=white,draw=black,line width=0.37pt] (6,-4) circle (0.3);
    \filldraw[fill=white,draw=black,line width=0.37pt] (7,-4) circle (0.3);
    \filldraw[fill=white,draw=black,line width=0.37pt] (8,-4) circle (0.3);
    \filldraw[fill=white,draw=black,line width=0.37pt] (9,-4) circle (0.3);
    \filldraw[fill=kgray,draw=black,line width=0.37pt] (10,-4) circle (0.3);
    \draw[black,line width=0.37pt] (9.788,-3.788) -- (10.212,-4.212) (9.788,-4.212) -- (10.212,-3.788);
    \filldraw[fill=kgray,draw=black,line width=0.37pt] (11,-4) circle (0.3);
    \draw[black,line width=0.37pt] (10.788,-3.788) -- (11.212,-4.212) (10.788,-4.212) -- (11.212,-3.788);
    \filldraw[fill=kgray,draw=black,line width=0.37pt] (12,-4) circle (0.3);
    \draw[black,line width=0.37pt] (11.788,-3.788) -- (12.212,-4.212) (11.788,-4.212) -- (12.212,-3.788);
    \filldraw[fill=kgray,draw=black,line width=0.37pt] (13,-4) circle (0.3);
    \draw[black,line width=0.37pt] (12.788,-3.788) -- (13.212,-4.212) (12.788,-4.212) -- (13.212,-3.788);
    \filldraw[fill=white,draw=black,line width=0.37pt] (14,-4) circle (0.3);
    \filldraw[fill=white,draw=black,line width=0.37pt] (15,-4) circle (0.3);
    \node[lab] at (1,-0.45) {1};
    \node[lab] at (2,-0.45) {2};
    \node[lab] at (3,-0.45) {3};
    \node[lab] at (4,-0.45) {4};
    \node[lab] at (5,-0.45) {5};
    \node[lab] at (6,-0.45) {6};
    \node[lab] at (7,-0.45) {7};
    \node[lab] at (8,-0.45) {8};
    \node[lab] at (9,-0.45) {9};
    \node[lab] at (10,-0.45) {10};
    \node[lab] at (11,-0.45) {11};
    \node[lab] at (12,-0.45) {12};
    \node[lab] at (13,-0.45) {13};
    \node[lab] at (14,-0.45) {14};
    \node[lab] at (15,-0.45) {15};
    \node[lab] at (0.45,-1) {1};
    \node[lab] at (0.45,-2) {2};
    \node[lab] at (0.45,-3) {3};
    \node[lab] at (0.45,-4) {4};
  \end{scope}
\end{tikzpicture}

\caption{The blackout of Proposition~\ref{prop:seas} for $h=4$ and $m=15$, with the 4 edges
of the complete bipartite graph between $\{1,2\}$ and $\{3,4\}$. Here $W_4+1=4$.}
\label{fig:seas}
\end{figure}

\begin{proof}[Proof of Proposition~\ref{thm:general}]
Theorem~\ref{thm:rowbound} gives the upper bound. Take $H$ to be the complete bipartite
graph with parts of sizes $\lfloor h/2\rfloor$ and $\lceil h/2\rceil$, so $q=q_h$, and
place the doubles at $p_j=1+(j-1)(W_h+1)$. The last double is in column
$1+(q_h-1)(W_h+1)=M_h$. Hence Proposition~\ref{prop:seas} gives a valid blackout with
$m+q_h$ points for every $m\ge M_h$.
\end{proof}

\begin{proof}[Proof of Proposition~\ref{prop:seas}]
Each double adds 1 point to a column, so $|Z|=m+q$. We check the 3 conditions of
Lemma~\ref{lem:criterion}.

\emph{Condition (i).} Only doubles have 2 black rows, and their pairs are distinct, so
$G_Z$ has no cycle of length $4$. A cycle $r_a\,c_x\,r_b\,c_y\,r_c\,c_z$ would need doubles
on $\{a,b\}$, $\{b,c\}$, and $\{a,c\}$, a triangle in $H$.

\emph{The black corners of a tilted rectangle.} Let $R$ be a tilted rectangle. In this
proof, \emph{the columns of $R$} are all columns from its smallest to its largest column,
not only the columns of its corners. The span of $R$ is at most $W_h$
(Lemma~\ref{lem:span}), so its columns contain at most 1 double. With no
double, all black points in the columns of $R$ lie in one boundary row, and $R$ has at most
1 black corner by (F1). Otherwise, shift the double to column $0$ and write its black rows as
$a<b$. Reflecting the rows if needed, we can assume that the seas in the columns of $R$ are
black in row $1$ to the left of column $0$ and in row $h$ to the right.

Let $K$ be the set of black corners of $R$. By (F1), the elements of $K$ can only be one
point $u=(x_1,1)$ with $x_1<0$, one point $v=(x_2,h)$ with $x_2>0$, and points $(0,c)$ with
$c\in\{a,b\}$. A corner $(0,1)$ excludes $u$, and a corner $(0,h)$ excludes $v$. We now
list the sets $K$ with $|K|\ge2$. In each case we find the other corners of $R$.
\begin{enumerate}[(a)]
\item $K=\{u,v\}$. These corners are opposite by (F2). The other diagonal has midpoint
$\bigl(\frac{x_1+x_2}2,\frac{1+h}2\bigr)$ and endpoints at this midpoint plus and minus a
vector $w$ with $|w|^2=\frac14\bigl((x_2-x_1)^2+(h-1)^2\bigr)$. By (F1) the endpoints lie in
rows $2$ to $h-1$. Hence $|w_2|\le\frac{h-3}2$ and $w_1^2>\frac14(x_2-x_1)^2$. So one
endpoint lies to the left of column $x_1$ and the other to the right of column $x_2$.
\item $K=\{(0,a),(0,b)\}$. These corners are opposite by (F3). The other diagonal has
midpoint $\bigl(0,\frac{a+b}2\bigr)$ and length $b-a$, and it is not vertical. Hence its
endpoints lie in columns $-w_1$ and $w_1$ for some $w_1\ne0$, and in rows strictly between
$a$ and $b$.
\item $K=\{u,d\}$ with $d=(0,c)$, and $u,d$ are opposite. Here $c\ge2$ by (F1). The other
diagonal has midpoint $\bigl(\frac{x_1}2,\frac{1+c}2\bigr)$ and endpoints at this midpoint
plus and minus a vector $w$ with $|w|^2=\frac14\bigl(x_1^2+(c-1)^2\bigr)$. The endpoints
are not in row $1$ by (F1). They are not in row $h$ by (F2), since the corner opposite to
$u$ is $d$. They are not in column $0$ by (F3), for the same reason. Since their rows are
at least $2$, we get $|w_2|\le\frac{c-3}2$. In particular $c\ge3$, and
$w_1^2>\frac14x_1^2$. Thus one endpoint
lies to the right of column $0$ and the other to the left of column $x_1$.
\item $K=\{u,d\}$ with $d=(0,c)$, and $u,d$ are adjacent. Again $c\ge2$. The side $d-u$ is
$(|x_1|,c-1)$, and the other corners are $u+n$ and $d+n$ with $n=t\,(-(c-1),|x_1|)$ for
some real $t\ne0$. The corner $u+n$ lies in row $1+t|x_1|$, so $t>0$. Hence both corners lie in
columns to the left of column $0$ and in rows at least $2$.
\item $K=\{v,d\}$ with $d=(0,c)$. The half turn $(x,j)\mapsto(-x,h+1-j)$ exchanges the 2
seas and fixes column $0$. It reduces this case to (c) or (d).
\item $|K|\ge3$. This is impossible. If $K$ contains $(0,a)$ and $(0,b)$, then the other 2
corners lie in rows strictly between $a$ and $b$ by (b), so they are neither $u$ nor $v$,
and they are not in column $0$. Otherwise $K=\{u,d,v\}$ with $d=(0,c)$. Then $u$ and $v$
are opposite by (F2) and $d$ is adjacent to both. So $(u-d)\cdot(v-d)=0$, that is
$x_1x_2=(c-1)(h-c)$. The left side is negative and the right side is not.
\end{enumerate}
In (a) to (d), the 2 other corners lie in columns different from $0$ and in rows $2$ to
$h-1$, or to the left of column $0$ and in rows at least $2$. The seas in these columns are
black only in row $1$ or only in row $h$, so the 2 other corners are white. By (e) the same
holds in the reflected cases. We have proved the following 3 statements for a tilted
rectangle $R$.
\begin{enumerate}[(1)]
\item $R$ has at most 2 black corners. Hence condition (ii) holds.
\item If $R$ has 2 opposite black corners, then its 2 white corners lie in rows $2$ to $h-1$,
one of them to the left and the other to the right of the double in the columns of $R$.
\item If $R$ has 2 adjacent black corners, then both white corners of $R$ lie on one side
of the double in the columns of $R$. On this side the seas in the columns of $R$ are black
in one boundary row, and the white corners are not in that row.
\end{enumerate}

\emph{Condition (iii).} Let $(R_1,R_2)$ be a mixed pair with common corners $p,q$ whose
support is black. A tilted rectangle of the pair has 2 black corners in the support, so its
corners $p,q$ are white by (1).

Suppose that $R_1$ is axis-aligned. Then $R_2$ is tilted with 2 opposite black corners. By
(2), $p$ and $q$ lie in different columns $x<z$ on the 2 sides of the double of $R_2$. So
the side $pq$ of $R_1$ is horizontal. The other 2 corners of $R_1$ are black points of the
columns $x$ and $z$ in a common row. But these columns are seas, one black only in row $1$
and the other only in row $h$, a contradiction.

Suppose that $R_2$ is axis-aligned. Then $R_1$ is tilted with 2 adjacent black corners. By
(3), $p=(x,j)$ and $q=(z,k)$ lie in sea columns that are black only in one boundary row
$\sigma$, and $j,k\ne\sigma$. The other corners of $R_2$ are $(x,k)$ and $(z,j)$, which are
white, a contradiction.

Suppose finally that both are tilted. The side opposite to $pq$ in $R_1$ is black, so (3)
applies to $R_1$. Let $D_1$ be the column of the double in the columns of $R_1$. Then $p$
and $q$ lie on one side of column $D_1$. The diagonal of $R_2$ other than $pq$ is black, so
(2) applies to $R_2$. Let $D_2$ be the column of the double in the columns of $R_2$. Then
$D_2$ lies strictly between the columns of $p$ and $q$. Hence $D_2\ne D_1$, and $D_2$ is a
second double in the columns of $R_1$, a contradiction. Thus (iii) holds, and $Z$ is valid by
Lemma~\ref{lem:criterion}.
\end{proof}

\begin{remark}\label{rem:spacing}
The proof uses the spacing $W_h+1$ only to keep 2 doubles out of the columns of one tilted
rectangle. The spacing $W_h$ is not always enough. Take the complete bipartite graph
between the rows $1,\dots,\lfloor h/2\rfloor$ and the other rows, with the doubles in
lexicographic order of their row pairs. Then the blackout with spacing $W_h$ is valid for
$h=3,5,7$ and not valid for $h=4,6,8$.
\end{remark}
